% ============================================================
%  Глава 1 — Математическая модель робота SAMURAI
% ============================================================
\chapter{Математическая модель робота SAMURAI}
\label{ch:model}

\section{Физическая постановка}
\label{sec:physical}

Робот SAMURAI --- гусеничное шасси с дифференциальным приводом.
Контроллер PCA9685 управляет двумя задними моторами (M1 левый, M2 правый);
передние моторы физически отсутствуют (см. \filepath{hardware\_presets/samurai\_default.yaml}).

\begin{filebox}[title={Параметры объекта (\filepath{config.yaml})}]
\begin{verbatim}
simulator:
  wheel_base:    0.17    # m  (track separation)
  max_linear:    0.30    # m/s
  max_angular:   2.00    # rad/s
\end{verbatim}
\end{filebox}

\textbf{Команда движения} формируется в виде вектора $\vect{u} = (u_v, u_\omega)\T$,
где $u_v$ --- желаемая линейная скорость [м/с], $u_\omega$ --- желаемая
угловая скорость [рад/с]. Дифференциальная кинематика преобразует $\vect{u}$
в скорости гусениц:
\begin{equation}
  v_L = u_v + \tfrac{B}{2}\,u_\omega, \qquad
  v_R = u_v - \tfrac{B}{2}\,u_\omega,
  \label{eq:diff_kinematics}
\end{equation}
где $B$ --- ширина базы (\texttt{wheel\_base}).

\section{Нелинейная ОДУ модели}
\label{sec:ode}

В отличие от тривиальной кинематической модели, реальный робот имеет
\textbf{инерционность моторов}: команда $u_v$ не отрабатывается мгновенно,
а через апериодическое звено первого порядка с постоянной времени $\tau_v$.
Аналогично для угловой скорости (постоянная $\tau_\omega$).

Введём \textbf{вектор состояния} ($n=5$):
\begin{equation}
  \vect{x}(t) = \big(\,p_x,\; p_y,\; \theta,\; v,\; \omega\,\big)\T
  \label{eq:state_vector}
\end{equation}

\begin{itemize}
  \item $p_x, p_y$ --- координаты центра робота в мировой системе [м];
  \item $\theta$ --- угол курса (yaw) [рад];
  \item $v$ --- фактическая линейная скорость [м/с];
  \item $\omega$ --- фактическая угловая скорость [рад/с].
\end{itemize}

\begin{formulabox}[title={Нелинейная модель в форме $\dot{\vect{x}} = \vect{f}(\vect{x}, \vect{u})$}]
\begin{equation}
  \begin{cases}
    \dot{p}_x      = v\cos\theta,\\[2pt]
    \dot{p}_y      = v\sin\theta,\\[2pt]
    \dot{\theta}   = \omega,\\[2pt]
    \dot{v}        = \dfrac{1}{\tau_v}\,(u_v - v),\\[6pt]
    \dot{\omega}   = \dfrac{1}{\tau_\omega}\,(u_\omega - \omega).
  \end{cases}
  \label{eq:nonlinear_ode}
\end{equation}
\end{formulabox}

Первые три уравнения --- кинематика дифференциального робота
(идентична уравнениям 1.1 из учебника Козлова, формулы для $\dot{x}$ и $\dot{z}$).
Последние два --- динамика моторов как RC-звено с постоянной времени $\tau$,
определяемой механической инерцией ротора и противо-ЭДС.

\begin{notebox}
Постоянные времени $\tau_v \approx 0{,}15$ с, $\tau_\omega \approx 0{,}10$ с
оценены экспериментально по графику разгона (см.
\filepath{tests/test\_motor\_lag.py}). Для других шасси значения нужно
переопределить --- они вынесены в \texttt{config.yaml} раздел \texttt{control.plant}.
\end{notebox}

\section{Линеаризация в опорной точке}
\label{sec:linearization}

Выберем \textbf{опорный режим} --- равномерное прямолинейное движение
вперёд со скоростью $v_0$:
\begin{equation}
  \bar{\vect{x}} = (\bar{p}_x,\; \bar{p}_y,\; 0,\; v_0,\; 0)\T,
  \qquad
  \bar{\vect{u}} = (v_0,\; 0)\T.
  \label{eq:operating_point}
\end{equation}

Подстановка $\bar{\vect{x}}, \bar{\vect{u}}$ в \eqref{eq:nonlinear_ode}
даёт правые части $\dot{\bar{\vect{x}}} = (v_0,\,0,\,0,\,0,\,0)\T$, что
соответствует нулевым отклонениям состояния от опорной траектории при
постоянном $v_0$.

Согласно методу линеаризации (см. учебник п.~1.1, формулы 1.2),
уравнение для отклонений $\Delta\vect{x} = \vect{x} - \bar{\vect{x}}$,
$\Delta\vect{u} = \vect{u} - \bar{\vect{u}}$ имеет вид:
\begin{equation}
  \Delta\dot{\vect{x}}(t) = \mat{A}\,\Delta\vect{x}(t) + \mat{B}\,\Delta\vect{u}(t),
  \qquad
  \mat{A} = \left.\frac{\partial \vect{f}}{\partial\vect{x}}\right|_{\bar{\vect{x}},\bar{\vect{u}}},
  \quad
  \mat{B} = \left.\frac{\partial \vect{f}}{\partial\vect{u}}\right|_{\bar{\vect{x}},\bar{\vect{u}}}.
  \label{eq:linearization}
\end{equation}

\subsection{Вычисление матриц Якоби}

Распишем частные производные функций $f_i$ по $x_j$:
\begin{align*}
  \frac{\partial f_1}{\partial \theta} &= -v\sin\theta\big|_0 = 0, &
  \frac{\partial f_1}{\partial v} &= \cos\theta\big|_0 = 1,\\
  \frac{\partial f_2}{\partial \theta} &= v\cos\theta\big|_0 = v_0, &
  \frac{\partial f_2}{\partial v} &= \sin\theta\big|_0 = 0,\\
  \frac{\partial f_3}{\partial \omega} &= 1, &
  \frac{\partial f_4}{\partial v} &= -1/\tau_v,\\
  \frac{\partial f_5}{\partial \omega} &= -1/\tau_\omega.
\end{align*}

Остальные элементы матрицы $\mat{A}$ равны нулю. Получаем:

\begin{formulabox}[title={Линеаризованная модель в пространстве состояний}]
\begin{equation}
  \mat{A} = \begin{pmatrix}
    0 & 0 & 0    & 1 & 0 \\
    0 & 0 & v_0  & 0 & 0 \\
    0 & 0 & 0    & 0 & 1 \\
    0 & 0 & 0    & -1/\tau_v & 0 \\
    0 & 0 & 0    & 0 & -1/\tau_\omega
  \end{pmatrix},
  \qquad
  \mat{B} = \begin{pmatrix}
    0 & 0 \\
    0 & 0 \\
    0 & 0 \\
    1/\tau_v & 0 \\
    0 & 1/\tau_\omega
  \end{pmatrix}.
  \label{eq:AB_matrices}
\end{equation}
\end{formulabox}

Матрица наблюдения $\mat{C}$ выбирается по доступным датчикам.
В роботе SAMURAI прямо измеряются $p_x, p_y$ (одометрия),
$\theta$ (IMU), а $v, \omega$ восстанавливаются из последовательных
позиций. Будем считать наблюдаемыми все 5 координат:
\begin{equation}
  \mat{C} = \mat{I}_5,
  \qquad
  \mat{D} = \mat{0}_{5\times 2}.
  \label{eq:CD_matrices}
\end{equation}

\subsection{Численные значения}
\label{sec:numerics}

При $v_0 = 0{,}2$ м/с (профиль \texttt{normal}), $\tau_v = 0{,}15$ с,
$\tau_\omega = 0{,}10$ с:
\begin{equation}
  \mat{A} = \begin{pmatrix}
    0 & 0 & 0   & 1 & 0\\
    0 & 0 & 0.2 & 0 & 0\\
    0 & 0 & 0   & 0 & 1\\
    0 & 0 & 0   & -6.667 & 0\\
    0 & 0 & 0   & 0 & -10
  \end{pmatrix},
  \quad
  \mat{B} = \begin{pmatrix}
    0 & 0\\
    0 & 0\\
    0 & 0\\
    6.667 & 0\\
    0 & 10
  \end{pmatrix}.
  \label{eq:AB_numerical}
\end{equation}

\section{Анализ управляемости}
\label{sec:controllability}

Объект полностью управляем тогда и только тогда, когда матрица управляемости
имеет полный ранг (учебник, формула перед п.~3.1):
\begin{equation}
  \mat{S}_y = \big[\,\mat{B}\,\big|\,\mat{A}\mat{B}\,\big|\,\mat{A}^2\mat{B}\,\big|\,\ldots\,\big|\,\mat{A}^{n-1}\mat{B}\,\big],
  \qquad
  \operatorname{rank} \mat{S}_y = n.
  \label{eq:controllability}
\end{equation}

Для нашей системы $n=5$, $r=2$, поэтому $\mat{S}_y$ имеет размер $5 \times 10$.
Вычисление в MATLAB (\filepath{matlab/analyze\_system.m}, см. главу~\ref{ch:implementation})
показывает $\operatorname{rank}\mat{S}_y = 5$ при $v_0 \neq 0$ ---
объект полностью управляем.

\begin{warnbox}
При $v_0 = 0$ опорная точка вырождена: координата $p_y$ становится
неуправляемой (нельзя двигаться вбок без вращения). Для синтеза регулятора
необходимо $v_0 > 0$ или модель в окрестности нетривиальной траектории.
\end{warnbox}

\section{Дискретизация}
\label{sec:discretization}

Для реализации регулятора в реальном времени переходим к дискретной модели
с периодом $T_s = 0{,}05$ с (20 Гц --- частота цикла \texttt{motor\_node}).
Согласно п.~1.5 учебника, дискретизация ZOH (zero-order hold) даёт:
\begin{equation}
  \boxed{
    \mat{A}_d = e^{\mat{A} T_s}, \qquad
    \mat{B}_d = \int_0^{T_s} e^{\mat{A}\tau}\,d\tau \cdot \mat{B},
  }
  \label{eq:zoh_discretization}
\end{equation}
что эквивалентно вычислению блочной матричной экспоненты:
\begin{equation}
  \exp\!\left(
    \begin{pmatrix} \mat{A} & \mat{B} \\ \mat{0} & \mat{0} \end{pmatrix} T_s
  \right)
  =
  \begin{pmatrix} \mat{A}_d & \mat{B}_d \\ \mat{0} & \mat{I} \end{pmatrix}.
  \label{eq:block_exp}
\end{equation}

Дискретная модель:
\begin{equation}
  \vect{x}_{k+1} = \mat{A}_d\,\vect{x}_k + \mat{B}_d\,\vect{u}_k,
  \qquad
  \vect{y}_k = \mat{C}\,\vect{x}_k.
  \label{eq:discrete_model}
\end{equation}

\subsection{Соответствие собственных значений}

Собственные значения непрерывной и дискретной систем связаны соотношением
(п.~3.1 учебника):
\begin{equation}
  \tilde{\lambda}_i = e^{\lambda_i T_s}.
  \label{eq:eigval_mapping}
\end{equation}
Условия устойчивости:
\begin{itemize}
  \item непрерывная: $\operatorname{Re}\lambda_i < 0$ ---
        собственные значения в левой полуплоскости;
  \item дискретная: $|\tilde{\lambda}_i| < 1$ ---
        внутри единичного круга.
\end{itemize}

\begin{notebox}
Для непрерывного робота $\mat{A}$ имеет три нулевых собственных значения
(апериодические интеграторы по $p_x, p_y, \theta$) и два отрицательных:
$-1/\tau_v$, $-1/\tau_\omega$. Объект \textbf{не устойчив асимптотически}
(но и не неустойчив --- это интегратор), что и ожидаемо: робот без
управления продолжает двигаться по инерции.
\end{notebox}
